% Options for packages loaded elsewhere
\PassOptionsToPackage{unicode}{hyperref}
\PassOptionsToPackage{hyphens}{url}
%
\documentclass[
  ignorenonframetext,
]{beamer}
\usepackage{pgfpages}
\setbeamertemplate{caption}[numbered]
\setbeamertemplate{caption label separator}{: }
\setbeamercolor{caption name}{fg=normal text.fg}
\beamertemplatenavigationsymbolsempty
% Prevent slide breaks in the middle of a paragraph
\widowpenalties 1 10000
\raggedbottom
\setbeamertemplate{part page}{
  \centering
  \begin{beamercolorbox}[sep=16pt,center]{part title}
    \usebeamerfont{part title}\insertpart\par
  \end{beamercolorbox}
}
\setbeamertemplate{section page}{
  \centering
  \begin{beamercolorbox}[sep=12pt,center]{part title}
    \usebeamerfont{section title}\insertsection\par
  \end{beamercolorbox}
}
\setbeamertemplate{subsection page}{
  \centering
  \begin{beamercolorbox}[sep=8pt,center]{part title}
    \usebeamerfont{subsection title}\insertsubsection\par
  \end{beamercolorbox}
}
\AtBeginPart{
  \frame{\partpage}
}
\AtBeginSection{
  \ifbibliography
  \else
    \frame{\sectionpage}
  \fi
}
\AtBeginSubsection{
  \frame{\subsectionpage}
}
\usepackage{amsmath,amssymb}
\usepackage{lmodern}
\usepackage{iftex}
\ifPDFTeX
  \usepackage[T1]{fontenc}
  \usepackage[utf8]{inputenc}
  \usepackage{textcomp} % provide euro and other symbols
\else % if luatex or xetex
  \usepackage{unicode-math}
  \defaultfontfeatures{Scale=MatchLowercase}
  \defaultfontfeatures[\rmfamily]{Ligatures=TeX,Scale=1}
\fi
% Use upquote if available, for straight quotes in verbatim environments
\IfFileExists{upquote.sty}{\usepackage{upquote}}{}
\IfFileExists{microtype.sty}{% use microtype if available
  \usepackage[]{microtype}
  \UseMicrotypeSet[protrusion]{basicmath} % disable protrusion for tt fonts
}{}
\makeatletter
\@ifundefined{KOMAClassName}{% if non-KOMA class
  \IfFileExists{parskip.sty}{%
    \usepackage{parskip}
  }{% else
    \setlength{\parindent}{0pt}
    \setlength{\parskip}{6pt plus 2pt minus 1pt}}
}{% if KOMA class
  \KOMAoptions{parskip=half}}
\makeatother
\usepackage{xcolor}
\IfFileExists{xurl.sty}{\usepackage{xurl}}{} % add URL line breaks if available
\IfFileExists{bookmark.sty}{\usepackage{bookmark}}{\usepackage{hyperref}}
\hypersetup{
  pdftitle={Regression Analysis Case Study - Hald Cement Dataset},
  pdfauthor={G Radhesyam},
  hidelinks,
  pdfcreator={LaTeX via pandoc}}
\urlstyle{same} % disable monospaced font for URLs
\newif\ifbibliography
\usepackage{longtable,booktabs,array}
\usepackage{calc} % for calculating minipage widths
\usepackage{caption}
% Make caption package work with longtable
\makeatletter
\def\fnum@table{\tablename~\thetable}
\makeatother
\usepackage{graphicx}
\makeatletter
\def\maxwidth{\ifdim\Gin@nat@width>\linewidth\linewidth\else\Gin@nat@width\fi}
\def\maxheight{\ifdim\Gin@nat@height>\textheight\textheight\else\Gin@nat@height\fi}
\makeatother
% Scale images if necessary, so that they will not overflow the page
% margins by default, and it is still possible to overwrite the defaults
% using explicit options in \includegraphics[width, height, ...]{}
\setkeys{Gin}{width=\maxwidth,height=\maxheight,keepaspectratio}
% Set default figure placement to htbp
\makeatletter
\def\fps@figure{htbp}
\makeatother
\setlength{\emergencystretch}{3em} % prevent overfull lines
\providecommand{\tightlist}{%
  \setlength{\itemsep}{0pt}\setlength{\parskip}{0pt}}
\setcounter{secnumdepth}{-\maxdimen} % remove section numbering
\ifLuaTeX
  \usepackage{selnolig}  % disable illegal ligatures
\fi

\title{Regression Analysis Case Study - Hald Cement Dataset}
\author{G Radhesyam}
\date{19 March 2022}

\begin{document}
\frame{\titlepage}

\begin{frame}{HALD CEMENT Dataset}
\protect\hypertarget{hald-cement-dataset}{}
\begin{itemize}
\item
  The response variable Y is the heat evolved in a cement mix.
\item
  The four explanatory variables are ingredients of the mix,

  \begin{itemize}
  \tightlist
  \item
    X1:Tricalcium aluminate
  \item
    X2:Tricalcium silicate,
  \item
    X3:Tetracalcium alumino ferrite
  \item
    X4:Dicalcium silicate.
  \end{itemize}
\end{itemize}
\end{frame}

\begin{frame}[fragile]{HALD CEMENT Dataset}
\protect\hypertarget{hald-cement-dataset-1}{}
\begin{verbatim}
## Rows: 13 Columns: 5
## -- Column specification --------------------------------------------------------
## Delimiter: ","
## dbl (5): Y, X1, X2, X3, X4
## 
## i Use `spec()` to retrieve the full column specification for this data.
## i Specify the column types or set `show_col_types = FALSE` to quiet this message.
\end{verbatim}

\begin{verbatim}
## # A tibble: 6 x 5
##       Y    X1    X2    X3    X4
##   <dbl> <dbl> <dbl> <dbl> <dbl>
## 1  78.5     7    26     6    60
## 2  74.3     1    29    15    52
## 3 104.     11    56     8    20
## 4  87.6    11    31     8    47
## 5  95.9     7    52     6    33
## 6 109.     11    55     9    22
\end{verbatim}

\begin{verbatim}
## spec_tbl_df [13 x 5] (S3: spec_tbl_df/tbl_df/tbl/data.frame)
##  $ Y : num [1:13] 78.5 74.3 104.3 87.6 95.9 ...
##  $ X1: num [1:13] 7 1 11 11 7 11 3 1 2 21 ...
##  $ X2: num [1:13] 26 29 56 31 52 55 71 31 54 47 ...
##  $ X3: num [1:13] 6 15 8 8 6 9 17 22 18 4 ...
##  $ X4: num [1:13] 60 52 20 47 33 22 6 44 22 26 ...
##  - attr(*, "spec")=
##   .. cols(
##   ..   Y = col_double(),
##   ..   X1 = col_double(),
##   ..   X2 = col_double(),
##   ..   X3 = col_double(),
##   ..   X4 = col_double()
##   .. )
##  - attr(*, "problems")=<externalptr>
\end{verbatim}

\begin{verbatim}
## cols(
##   Y = col_double(),
##   X1 = col_double(),
##   X2 = col_double(),
##   X3 = col_double(),
##   X4 = col_double()
## )
\end{verbatim}
\end{frame}

\begin{frame}{HALD CEMENT - FINAL MODEL SUMMARY}
\protect\hypertarget{hald-cement---final-model-summary}{}
\begin{itemize}
\item
  Model : X1+X2

  \begin{itemize}
  \tightlist
  \item
    Y = 52.5773 + (1.4683) X1 + ( 0.6623) X2
  \item
    All Regression assumptions met.
  \end{itemize}
\item
  Normality test

  \begin{itemize}
  \tightlist
  \item
    Shapiro-wilks Test - Distribution is Normal
  \end{itemize}
\item
  W = 0.90219, p-value = 0.1433
\item
  Anderson-Darling test - Distribution is normal
\item
  A = 0.61361, p-value = 0.08628
\item
  Equal Variance Test

  \begin{itemize}
  \tightlist
  \item
    studentized Breusch-Pagan test
  \item
    BP = 0.16688, df = 2, p-value = 0.9199
  \item
    H0 : Variances are equal
  \end{itemize}
\end{itemize}
\end{frame}

\begin{frame}[fragile]{HALD CEMENT - FINAL MODEL SUMMARY (Contd\ldots)}
\protect\hypertarget{hald-cement---final-model-summary-contd}{}
\begin{verbatim}
  * Independence/Auto Correllation Test (Durbin-Watson (D-W) test )
     - H0 : Auto correlation is zero
\end{verbatim}

\begin{longtable}[]{@{}llll@{}}
\toprule
lag & Autocorrelation & D-W Statistic & p-value \\
\midrule
\endhead
1 & -0.05450402 & 1.92164 & 0.892 \\
\bottomrule
\end{longtable}

\begin{itemize}
\tightlist
\item
  Linearity test

  \begin{itemize}
  \item
    Box-Tidwell - H0 : No Transformation needed
  \item
    Box-Tidwell test (Y,X1)
  \item
    MLE of lambda Score Statistic (z) Pr(\textgreater\textbar z\textbar)
  \item
\begin{verbatim}
  0.23145             -0.7366   0.4614
\end{verbatim}
  \item
    Box-Tidwell test ( Y,X2)
  \item
    MLE of lambda Score Statistic (z) Pr(\textgreater\textbar z\textbar)
  \item
    -0.71695 -1.35 0.177
  \end{itemize}
\end{itemize}
\end{frame}

\begin{frame}[fragile]{HALD CEMENT - MODEL ANALYSIS and DETAILS}
\protect\hypertarget{hald-cement---model-analysis-and-details}{}
\begin{itemize}
\item
  An F statistic is a value you get when you run an ANOVA test or a
  regression analysis to find out if the means between two populations
  are significantly different.
\item
  It's similar to a T statistic from a T-Test;
\item
  A T-test will tell you if a single variable is statistically
  significant and
\item
  an F test will tell you if a group of variables are jointly
  significant.
\item
  \url{https://www.statisticshowto.com/probability-and-statistics/f-statistic-value-test/}
\item
  Regression NULL Hypothesis: All of the regression coefficients are
  equal to zero
\item
  Tip: P-value is High ,NUll will Fly
\item
\begin{verbatim}
 P-value is low, NULL will blow
\end{verbatim}
\item
  T-Test in Regression
\item
  \url{https://www.statology.org/t-test-linear-regression/}
\end{itemize}

\begin{verbatim}
## 
## Call:
## lm(formula = haldcement$Y ~ ., data = haldcement)
## 
## Residuals:
##     Min      1Q  Median      3Q     Max 
## -3.1750 -1.6709  0.2508  1.3783  3.9254 
## 
## Coefficients:
##             Estimate Std. Error t value Pr(>|t|)  
## (Intercept)  62.4054    70.0710   0.891   0.3991  
## X1            1.5511     0.7448   2.083   0.0708 .
## X2            0.5102     0.7238   0.705   0.5009  
## X3            0.1019     0.7547   0.135   0.8959  
## X4           -0.1441     0.7091  -0.203   0.8441  
## ---
## Signif. codes:  0 '***' 0.001 '**' 0.01 '*' 0.05 '.' 0.1 ' ' 1
## 
## Residual standard error: 2.446 on 8 degrees of freedom
## Multiple R-squared:  0.9824, Adjusted R-squared:  0.9736 
## F-statistic: 111.5 on 4 and 8 DF,  p-value: 4.756e-07
\end{verbatim}
\end{frame}

\begin{frame}[fragile]{HALD CEMENT - Scatter plot}
\protect\hypertarget{hald-cement---scatter-plot}{}
\includegraphics{E:/Radhesyam/RProjects/MDSC206/Regression/2022/RegressionAnalysisCaseStudy_2022_files/figure-beamer/unnamed-chunk-4-1.pdf}

\begin{itemize}
\tightlist
\item
  \url{https://ggobi.github.io/ggally/reference/ggpairs.html}
\item
\end{itemize}

\begin{verbatim}
## Registered S3 method overwritten by 'GGally':
##   method from   
##   +.gg   ggplot2
\end{verbatim}

\includegraphics{E:/Radhesyam/RProjects/MDSC206/Regression/2022/RegressionAnalysisCaseStudy_2022_files/figure-beamer/unnamed-chunk-5-1.pdf}
\includegraphics{E:/Radhesyam/RProjects/MDSC206/Regression/2022/RegressionAnalysisCaseStudy_2022_files/figure-beamer/unnamed-chunk-6-1.pdf}
\end{frame}

\begin{frame}[fragile]{HALD CEMENT - Influence points}
\protect\hypertarget{hald-cement---influence-points}{}
\begin{itemize}
\tightlist
\item
  \url{https://www.rdocumentation.org/packages/stats/versions/3.6.2/topics/influence.measures}
\item
  influence.measures: Regression Deletion Diagnostics
\item
  \url{https://rpubs.com/mpfoley73/501093} \tiny
\end{itemize}

\begin{verbatim}
##  1  3  5 10 
##  1  3  5 10
\end{verbatim}

\begin{verbatim}
## Potentially influential observations of
##   lm(formula = haldcement$Y ~ ., data = haldcement) :
## 
##    dfb.1_  dfb.X1  dfb.X2  dfb.X3  dfb.X4  dffit cov.r   cook.d hat  
## 1   0.00    0.00    0.00    0.00    0.00    0.00  4.34_*  0.00   0.55
## 3  -1.10_*  1.03_*  1.08_*  1.08_*  1.10_* -1.24  2.19    0.30   0.58
## 5  -0.02    0.00    0.02    0.00    0.02    0.09  3.00_*  0.00   0.36
## 10 -0.05    0.12    0.04    0.08    0.05    0.30  6.33_*  0.02   0.70
\end{verbatim}

\begin{verbatim}
## [1] 4 9
\end{verbatim}

\normalsize
\end{frame}

\begin{frame}{Influence points analysis}
\protect\hypertarget{influence-points-analysis}{}
\begin{itemize}
\item
  n= 4 and k= 9
\item
  A Point is an influence point if

  \begin{itemize}
  \item
    COOK's D-statistic(D) \textgreater{} 1
  \item
    DFBETAS \textgreater{}
    \({2 \over \sqrt{n}} = {2 \over \sqrt{4}} = 1\), i.e DFBETAS
    \textgreater{} 1
  \item
    DFFITS \textgreater{}
    \(2 * \sqrt{ k \over n} = 2 * \sqrt{9 \over 4} = 3\) , i.e DFFITS
    \textgreater{} 3
  \end{itemize}
\item
  INFLUENCE POINT
\item
  Point 3 (as per DFBETAS). None as per DFFITS and cook's distance D.
\item
  INVESTIGATION
\item
  Point 3 is within 2 sigma limits of X1,X2,X3 and X4. So retain it
\end{itemize}
\end{frame}

\begin{frame}[fragile]{MULTICOLLINEARITY Check}
\protect\hypertarget{multicollinearity-check}{}
\begin{verbatim}
##        X1        X2        X3        X4 
##  38.49621 254.42317  46.86839 282.51286
\end{verbatim}

\begin{itemize}
\tightlist
\item
  OBSERVATION: SEVERE MULTICOLLINEARITY as VIF \textgreater{} 10 FOR ALL
  VARIABLES
\end{itemize}
\end{frame}

\begin{frame}[fragile]{Correlation Co-efficients - Pearson test}
\protect\hypertarget{correlation-co-efficients---pearson-test}{}
\tiny

\begin{verbatim}
##    row column        cor            p
## 1    Y     X1  0.7307175 4.552045e-03
## 2    Y     X2  0.8162526 6.648249e-04
## 3   X1     X2  0.2285795 4.525662e-01
## 4    Y     X3 -0.5346707 5.976232e-02
## 5   X1     X3 -0.8241338 5.308752e-04
## 6   X2     X3 -0.1392424 6.500589e-01
## 7    Y     X4 -0.8213050 5.762318e-04
## 8   X1     X4 -0.2454451 4.189332e-01
## 9   X2     X4 -0.9729550 2.404645e-08
## 10  X3     X4  0.0295370 9.236906e-01
\end{verbatim}

\normalsize

\begin{itemize}
\tightlist
\item
  Test of Hypothesis (Pearson)

  \begin{itemize}
  \tightlist
  \item
    H0: \(\rho\) = 0 ; i.e No relationship between the variables.
  \item
    H1 : \(\rho\) \(\neq\) 0
  \end{itemize}
\item
  OBSERVATION:
\item
  No Correllation : Accept H0 : X1,X2 (p=0.4); X2,X3
  (p=0.6),X1,X4(p=0.4),X3,X4 (p=0.9)
\item
  Correlation exists : Accept H1 : X1,X3 (\textless{} 0.05);X2,X4
  (\textless{} 0.05)
\end{itemize}
\end{frame}

\begin{frame}[fragile]{Variable Selection - Forward selection}
\protect\hypertarget{variable-selection---forward-selection}{}
\begin{verbatim}
##   (Intercept) X1 X2 X3 X4        R2     SSres     AdjR2         Cp        BIC
## 1           1  0  0  0  1 0.6745420 883.86692 0.6449549 138.730833  -9.462884
## 2           1  1  0  0  1 0.9724710  74.76211 0.9669653   5.495851 -39.007873
## 3           1  1  1  0  1 0.9823355  47.97273 0.9764473   3.018233 -42.210745
## 4           1  1  1  1  1 0.9823756  47.86364 0.9735634   5.000000 -39.675391
\end{verbatim}

\begin{itemize}
\tightlist
\item
  RECOMMENDATION : X4 +X1 =97\%
\end{itemize}
\end{frame}

\begin{frame}[fragile]{Variable Selection - Backward Elimination}
\protect\hypertarget{variable-selection---backward-elimination}{}
\begin{verbatim}
##   (Intercept) X1 X2 X3 X4        R2     SSres     AdjR2         Cp        BIC
## 1           1  0  1  0  0 0.6662683 906.33634 0.6359290 142.486407  -9.136532
## 2           1  1  1  0  0 0.9786784  57.90448 0.9744140   2.678242 -42.329587
## 3           1  1  1  0  1 0.9823355  47.97273 0.9764473   3.018233 -42.210745
## 4           1  1  1  1  1 0.9823756  47.86364 0.9735634   5.000000 -39.675391
\end{verbatim}

\begin{itemize}
\tightlist
\item
  RECOMMENDATION : X2 + X1 = 97\%
\end{itemize}
\end{frame}

\begin{frame}[fragile]{Variable Selection - Stepwise Regression}
\protect\hypertarget{variable-selection---stepwise-regression}{}
\begin{verbatim}
##   (Intercept) X1 X2 X3 X4        R2     SSres     AdjR2         Cp        BIC
## 1           1  0  0  0  1 0.6745420 883.86692 0.6449549 138.730833  -9.462884
## 2           1  1  1  0  0 0.9786784  57.90448 0.9744140   2.678242 -42.329587
## 3           1  1  1  1  0 0.9822847  48.11061 0.9763796   3.041280 -42.173433
## 4           1  1  1  1  1 0.9823756  47.86364 0.9735634   5.000000 -39.675391
\end{verbatim}

\begin{itemize}
\tightlist
\item
  RECOMMENDATION : X4+X1+X2= 97\%
\end{itemize}
\end{frame}

\begin{frame}{Variable Selection - PCA}
\protect\hypertarget{variable-selection---pca}{}
\includegraphics{E:/Radhesyam/RProjects/MDSC206/Regression/2022/RegressionAnalysisCaseStudy_2022_files/figure-beamer/unnamed-chunk-19-1.pdf}
* RECOMMENDATION : : PCA1: X2 and X4 ( X4 very close to X3)
\end{frame}

\begin{frame}{FINALIZING VARIABLE/MODEL selection}
\protect\hypertarget{finalizing-variablemodel-selection}{}
\begin{longtable}[]{@{}lllll@{}}
\toprule
VAR & PRESS & VIF & \(R^{2}\)(adj) & \(R^{2}\)(pred) \\
\midrule
\endhead
X1+X2 & 93.8 & 1 & 97\% & 96\% \\
X2+X4 & 1461.8 & 18 & 61\% & 46\% \\
X4+X1 & 121.2 & 1 & 96\% & 95\% \\
X2+X3 & 701.7 & 1 & 81\% & 74\% \\
\bottomrule
\end{longtable}

\begin{itemize}
\tightlist
\item
  RECOMMENDATION :PRESS stat should be small ( so that R-sq(pred) is
  high) and VIF=1.
\item
  We will go with Model 1 : X1+X2
\end{itemize}
\end{frame}

\begin{frame}{FINAL Model : X1 + X2 : Validation of Regression
Assumptions}
\protect\hypertarget{final-model-x1-x2-validation-of-regression-assumptions}{}
\begin{itemize}
\item
  Normality test

  \begin{itemize}
  \tightlist
  \item
    Shapiro-wilks Test - Distribution is Normal
  \end{itemize}
\item
  W = 0.90219, p-value = 0.1433
\item
  Anderson-Darling test - Distribution is normal
\item
  A = 0.61361, p-value = 0.08628
\item
  Equal Variance Test

  \begin{itemize}
  \tightlist
  \item
    studentized Breusch-Pagan test
  \item
    BP = 0.16688, df = 2, p-value = 0.9199
  \end{itemize}
\end{itemize}
\end{frame}

\begin{frame}[fragile]{FINAL Model : X1 + X2 : Validation of Regression
Assumptions (contd..)}
\protect\hypertarget{final-model-x1-x2-validation-of-regression-assumptions-contd..}{}
\begin{itemize}
\tightlist
\item
  Independence/Auto Correllation Test (Durbin-Watson (D-W) test )
\end{itemize}

\begin{longtable}[]{@{}llll@{}}
\toprule
lag & Autocorrelation & D-W Statistic & p-value \\
\midrule
\endhead
1 & -0.05450402 & 1.92164 & 0.892 \\
\bottomrule
\end{longtable}

\begin{itemize}
\tightlist
\item
  Linearity test

  \begin{itemize}
  \item
    Box-Tidwell test (Y,X1)
  \item
    MLE of lambda Score Statistic (z) Pr(\textgreater\textbar z\textbar)
  \item
\begin{verbatim}
  0.23145             -0.7366   0.4614
\end{verbatim}
  \item
    Box-Tidwell test ( Y,X2)
  \item
    MLE of lambda Score Statistic (z) Pr(\textgreater\textbar z\textbar)
  \item
    -0.71695 -1.35 0.177
  \end{itemize}
\end{itemize}
\end{frame}

\end{document}
